\section{파생 기법}\label{sec:derived}

순수 ASWISS의 약점은 둘이다.
\begin{enumerate}
  \item 정보집합별 표 없이 계보를 따르는 비용이 분할 단계에 대해 곱으로 늘어난다(\ref{sec:exp-route}절).
  \item 라벨에 필요한 반복 전략을 모두 저장할 수 없어 표본으로 근사한다.
\end{enumerate}
이 절은 이 약점을 고치는 기법들을 정리한다. 각 기법이 정리의 가정을 얼마나 지키는지 함께 적는다.

\subsection{평균 전략 방향}\label{sec:derived-avg}
epoch의 반복 전략 전부 대신 \emph{평균 전략 하나}로 라벨을 구하면 저장과 계산이 epoch 길이만큼 줄어든다.
두 방향은 일반적으로 같지 않다. 그 차이를 정확히 나타낸 것이 다음 정리다.

\paragraph{기호.}
플레이어 $p$의 정보집합 $I$, 행동 $a$, 프로필 열 $\sigma_1,\dots,\sigma_T$를 고정한다. $\bar\sigma$는 식~\eqref{eq:avg}의 평균이다.
\begin{itemize}
  \item 상대와 우연의 실현 계획: $y_t(z)=\pi^{\sigma_t}_{-p}(z)$, $\bar y=\frac1T\sum_ty_t$
  \item $Ia$에서 $z$까지 자기 행동 확률: $w_t^{Ia}(z)=\pi_p^{\sigma_t}(Ia\to z)$, $\bar w^{Ia}=\frac1T\sum_tw_t^{Ia}$
  \item $\hat\sigma_{Ia}$: $Ia$ 아래 $p$의 결정만 $Ia$ 조건부 평균으로 바꾼 $\bar\sigma$
  \item 도달확률 가중 항: $\Omega_T(I,a)=v^{\hat\sigma_{Ia}}(I,a)-v^{\bar\sigma}(I,a)$
  \item 공분산 항: $K_T(I,a)=\frac1T\sum_t\sum_{z\sqsupseteq Ia}(y_t-\bar y)(z)\,(w_t^{Ia}-\bar w^{Ia})(z)\,u_p(z)$
  \item 행동과 무관한 상수: $c_T(I)=T\langle\bar\sigma(I),v^{\bar\sigma}(I)\rangle-\sum_t\langle\sigma_t(I),v^{\sigma_t}(I)\rangle$
\end{itemize}
\begin{theorem}[평균 전략 방향의 분해]\label{thm:avgdir}
모든 행동 $a$에 대해 다음이 성립한다.
\[
\sum_{t=1}^Tr^{\sigma_t}(I,a)=T\,r^{\bar\sigma}(I,a)+T\,\Omega_T(I,a)+T\,K_T(I,a)+c_T(I).
\]
\end{theorem}
\begin{itemize}
  \item 상대의 전략이 어떻게 바뀌어도 차이를 만들지 않는다. 차이는 자기 후속 전략의 도달확률 가중($\Omega$)과, 상대와 자기 후속 전략의 시간 공분산($K$)뿐이다.
  \item $I$ 아래에 $p$의 결정이 없으면(전선 정보집합) $\Omega=K=0$이고, 행동 순위가 정확히 같다.
  \item CFR의 반복 전략은 진동할 수 있으므로 $K\to0$의 보장은 없다.
  \item 그래서 정리~\ref{thm:epoch}, \ref{thm:replay}의 보장은 평균 전략 방향으로 옮겨지지 않는다.
\end{itemize}
\paragraph{실험(mini-stud, 외부 샘플링, 학습 1,800초, seed 1).}
분할 시점 전략 하나로 방향을 구해 계보를 따르는 판(원안)에서 그 전략만 바꿔 보았다(2.62억 반복).
\begin{itemize}
  \item 순간 전략 하나를 평균 전략으로 바꾸자 0.147에서 0.035로, 같은 반복의 exploitability가 4배 낮아졌다.
  \item 분할 판정까지 평균 전략의 정확한 방향으로 하면 0.012였다. 같은 반복의 창 라벨(0.041)보다 3.5배 낮고 bucket은 28\% 적었다.
\end{itemize}

\subsection{bucket 이력 시작 칸}\label{sec:derived-history}
\begin{itemize}
  \item \textbf{문제.} 합법 행동집합만으로 묶은 시작 분할은 한 경로가 같은 칸을 두 번 지날 수 있다(absent-mindedness). 예를 들어 같은 플레이어의 5구와 7구 결정, 또는 레이즈 싸움의 반복된 결정이 같은 칸에 든다.
  \item \textbf{정의.} 그래서 시작 칸을 "그 플레이어가 이번 판에서 지나온 자기 칸(constraint)의 순서" $(i,c_1,\dots,c_k)$로 정한다. 그러면 한 경로가 같은 칸을 두 번 지나지 않는다.
  \item \textbf{칸 수.} mini-stud에서는 20칸이다. 헤즈업 세븐 포커에서는 constraint를 (스트리트, 첫/둘째 행동자, 합법 행동집합)으로 두면 규칙으로 세어 284칸이다.
  \item \textbf{정리와의 관계.} 5절의 보장은 absent-mindedness를 금지하지 않아도 성립한다. 원 게임에서 lifting된 전략은 여전히 행동 전략이기 때문이다. 이 기법의 목적은 추상 게임 학습의 질이다.
\end{itemize}

\subsection{관측 분할과 전략 분할의 이중 구조}\label{sec:derived-pp}
계보 라우팅의 비용을 없애려고, 정보집합이 들어갈 칸을 정하는 일과 전략표를 나누는 일을 분리한다.
설계 문서에서는 이것을 prior–posterior 구조라고 불렀다. 이 논문에서는 하는 일을 따라 \emph{관측 분할}(observation partition, 옛 이름 prior)과 \emph{전략 분할}(strategy partition, 옛 이름 posterior)이라 부른다.
\begin{itemize}
  \item \textbf{관측 분할.}
    \begin{itemize}
      \item 정보집합의 관측(카드, 베팅 등)의 특징과 문턱값으로 만든 결정 트리다.
      \item 잎(모둠)마다 전략표 하나를 가리키는 포인터가 있다. 여러 잎이 한 표를 같이 쓸 수 있다.
      \item 정보집합의 칸은 트리를 한 번 내려가면 정해지므로, 계보를 다시 계산할 필요가 없다.
    \end{itemize}
  \item \textbf{관측 분할의 세분.}
    \begin{itemize}
      \item 잎마다 후보 자르기(특징, 문턱값)를 몇 개 들고, 각 쪽의 후회 합을 쌓는다.
      \item 양쪽 방문이 합쳐 1,000이 되면 한 번 시험한다.
      \item 통과하면 잎을 둘로 나누고, 두 자식은 같은 표를 계속 가리킨다(표는 그대로).
    \end{itemize}
  \item \textbf{전략 분할의 세분.}
    \begin{itemize}
      \item 정해진 간격마다 각 표를 쓰는 잎들을 기록의 argmax 행동으로 묶는다.
      \item 무리마다 "그 무리 대 나머지"를 시험한다.
      \item 통과한 무리는 새 표로 옮긴다.
    \end{itemize}
\end{itemize}
두 시험은 모두 \emph{두 묶음의 숨은 후회}를 쓴다.
\[
\Delta=\psi(R_1)+\psi(R_2)-\psi(R_1+R_2)\ \ge0.
\]
이것은 보조정리~\ref{lem:label}의 숨은 후회를 두 묶음으로 나눈 것이다. 잡음 항은 둘 중 하나를 쓴다.
\begin{itemize}
  \item 정규 근사: $\Delta\ge\varepsilon t+z_\delta(\hat\sigma_1\sqrt{n_1}+\hat\sigma_2\sqrt{n_2})$
  \item Freedman 부등식~\citep{freedman1975tail}의 해: $x_j=c+\sqrt{c^2+2V_j\ln(1/\delta)}$, $c=L\ln(1/\delta)/3$로 두고 $\Delta\ge\varepsilon T+x_1+x_2$
\end{itemize}
$\varepsilon$은 "나눠서 숨은 후회가 반복당 이만큼은 줄어야 나눈다"는 문턱이다.

\paragraph{장기 실행(mini-stud, seed 1).}
\begin{itemize}
  \item 설정: bucket 이력 시작 칸, 관측 분할 실시간 시험, 전략 분할 실시간 Freedman 시험, 새 표는 부모 표를 복사(완전 상속)
  \item 이 설정으로 $\varepsilon=0,0.01,0.03,0.1$을 167.77억 반복까지 돌렸다(표~\ref{tab:longrun}).
\end{itemize}
\begin{table}[h]
\centering\small
\caption{mini-stud 장기 실행, 평균 정책 exploitability(괄호는 쓰이는 표 수)}\label{tab:longrun}
\begin{tabular}{@{}lrrrr@{}}
\toprule
반복 & $\varepsilon=0$ & $\varepsilon=0.01$ & $\varepsilon=0.03$ & $\varepsilon=0.1$ \\
\midrule
1.31억 & 0.528 (5,668) & 0.552 (613) & 0.639 (250) & 0.571 (68) \\
20.97억 & 0.338 (14,790) & 0.506 (871) & 0.623 (311) & 0.626 (77) \\
167.77억 & \textbf{0.173} (20,353) & 0.496 (962) & 0.661 (340) & 0.592 (82) \\
\bottomrule
\end{tabular}
\end{table}
\begin{itemize}
  \item \textbf{$\varepsilon=0$.} 끝까지 내려간 것은 $\varepsilon=0$뿐이다. 숨은 후회도 반복당 2에서 0.066으로 줄었다.
  \item \textbf{$\varepsilon>0$.} 10억 반복 근처부터 평평했고, 숨은 후회가 3–4에서 줄지 않았다.
  \item \textbf{k-means 기준선.} 같은 학습기의 k-means 고정 추상화는 900초(약 2.24억 반복)까지만 있고, 1.31억 반복에서 0.339였다. 그 구간에서는 이 판이 더 높았다. 그 뒤의 비교는 기준선이 없어서 하지 않았다.
  \item \textbf{같은 반복의 다른 판.} 같은 반복에서는 이 판이 순수판(표~\ref{tab:pure}, 1.31억 반복 0.0140)보다 훨씬 높다.
    \begin{itemize}
      \item 표를 완전 상속한 것과 관측 분할의 칸이 다른 것이 섞여 있어서 원인을 나누지 않았다.
      \item 이 구조의 이득은 계보 비용이 없다는 데 있다.
    \end{itemize}
\end{itemize}

\subsection{고정 문턱의 $\varepsilon$-수렴과 1단 게임의 반례}\label{sec:derived-theta}
분할을 "추정한 숨은 후회 $\hat C_W(B)$가 $\theta\abs W$를 넘을 때"로 정하는 판을 생각한다.
\begin{itemize}
  \item \textbf{주장.} 고정 문턱 $\theta>\eta>0$으로 판정하면 다음이 성립한다는 개요가 있다.
    \begin{itemize}
      \item 분할은 유한 번이다.
      \item 그 뒤 모든 bucket의 숨은 후회 비율이 $\theta+\eta$ 아래에 있다.
      \item $\NC\le O(1/\sqrt{\abs W})+\abs{\cB^*}(\theta+\eta)$이다(\texttt{ASWISS\_PROOF\_OUTLINE\_KO.md}).
    \end{itemize}
  \item \textbf{논리의 상태.} 논리는 정리~\ref{thm:epoch}와 같은 뼈대다. 다만 버킷 크기와 무관한 표본 추정(그 개요의 L4)은 아직 증명하지 않았다. 라벨 항등식으로 선형화하는 방법을 제안해 두었다.
  \item \textbf{우리 실행과의 대응.} \ref{sec:derived-pp}절의 $\varepsilon>0$ 실행이 이 문턱판에 해당한다. 평평해진 결과와 결이 맞는다.
\end{itemize}
같은 개요가 쓴 1단 행렬 게임 실험(공개 우연이 copy를 고르고 두 플레이어가 copy를 봄)을 다시 돌렸다. 기록된 수치가 모두 재현되었다.
\begin{itemize}
  \item \textbf{반례 B.} 라벨대로 한 플레이어만 나누면 NashConv가 0.238에서 2.5로 오른다. 다음 분할에서 0이 된다.
    \begin{itemize}
      \item 분할 전 묶인 추상 균형은 $x=y=(10/21,11/21)$이다.
      \item 분할 뒤 P1의 추상 maximin은 하나로 정해진 순수에 가까운 전략이다.
      \item 그래서 묶인 상대가 원 게임에서 크게 공략한다.
    \end{itemize}
  \item \textbf{$\varepsilon$-계열.} 분할 이득이 0.00023밖에 안 되는 분할도 2.5로 오를 수 있다. 문턱이 있으면 이런 분할을 하지 않는다.
  \item \textbf{분할 규칙 비교.} 같은 게임 2,892개에서, 한 번이라도 상한이 오른 게임의 비율은 다음과 같았다.
    \begin{itemize}
      \item 한 단계에 한 bucket을 나누는 규칙: 21.2\%
      \item 순수 ASWISS처럼 라벨이 갈린 bucket을 모두 동시에 나누는 규칙: 10.1\%
    \end{itemize}
\end{itemize}
이 게임에는 사적 정보가 없고, 균형이 하나인 경우만 썼다. 그래서 포커로 그대로 옮기지 않는다.

\subsection{헤즈업 세븐 포커}\label{sec:derived-7stud}
실제 크기의 헤즈업 세븐 포커(5–7구 베팅, 4구 버림과 공개는 고정 휴리스틱)에서는 순수판을 돌릴 수 없다. 원 정보집합이 학습 중 거의 다시 오지 않고, 계보 비용도 감당할 수 없다.
이 게임의 ASWISS(\texttt{hist\_bbf})는 다음을 함께 썼다.
\begin{itemize}
  \item 평균 전략 방향(\ref{sec:derived-avg}절)
  \item 정보집합별 표 대신 bucket 이력 표
  \item 역방향 분할
  \item 카드 군집 codebook을 관측으로 쓰는 사전 정보
  \item 학습기 MCCFR+
\end{itemize}
평가는 Policy-LBR(국소 최선 응답 공격자가 따는 양, ante/판)이다. exploitability의 근사 하한이지 exploitability가 아니다.
\begin{table}[h]
\centering\small
\caption{세븐 포커, 같은 학습 CPU 24시간(스레드 하나, seed 1, 평균 정책, 64파티클, 같은 20,000쌍 deal, 95\%)}\label{tab:7stud}
\resizebox{\linewidth}{!}{%
\begin{tabular}{@{}lrrr@{}}
\toprule
비교군 & 학습량 & LBR & k-means 256과의 짝지은 차이 \\
\midrule
ASWISS (파생 기법 + codebook) & 1.28억 roots, bucket 1,199,549 & \textbf{0.995 $\pm$ 0.116} & $-0.548\pm0.170$ \\
k-means 256 & 2.56억 roots, bucket 37,277 & 1.543 $\pm$ 0.121 & 0 \\
MMD & 239.5만 판 & 27.10 $\pm$ 1.77 & $+25.55\pm1.77$ \\
ESCHER & 53.1만 판 & 32.28 $\pm$ 1.80 & $+30.73\pm1.81$ \\
ReBeL & 23.9만 판 & 2.439 $\pm$ 0.246 & $+0.896\pm0.270$ \\
\bottomrule
\end{tabular}}
\end{table}
\begin{itemize}
  \item \textbf{같은 roots에서의 비율.} ASWISS/k-means 256 비율은 0.15M roots에서 1.50, 10.23M에서 0.82, 81.91M에서 0.70이었다. 초반에는 ASWISS가 나쁘고 1천만 roots 무렵부터 앞섰다.
  \item \textbf{이득의 출처.} ASWISS와 k-means 256은 같은 자료로 만든 codebook을 공유하고, ASWISS의 bucket 수가 32배다. 그래서 이득이 방법에서 왔는지 크기에서 왔는지는 이 실험으로 나눌 수 없다.
  \item \textbf{숨은 후회 판.} 관측·전략 이중 구조와 숨은 후회 시험의 세븐 포커 판(\texttt{pphist\_freedman\_rt})은 같은 조건의 비교군으로 넣어 두었고, 결과는 아직 없다.
\end{itemize}
